\section{结论}\label{sec:conclusion}
本文研究同步单天线站对未知宽带强弱重叠源的直接定位，以最小二乘估计消去波形系数，将问题化为位置优化。多候选联合AP保留少量粗搜索候选，分别联合调整源位置，再按最终残差选择结果；波形消元块NOMP在每次局部位置试探时重新估计系数，减少固定系数所造成的重复小步。

2400条观测的比较表明，多候选与联合精修分别贡献了可靠性和效率。五候选MCJ-AP达到99\%的双源20\,m成功率，相对单候选恢复96条，平均耗时0.718\,s。相同五候选的完整配对中，联合精修保持相同成功结果，将平均耗时从1.084\,s降至0.744\,s，减少31.36\%，并减少约42.30\%的梯度计算。共同240条上，MCJ-AP以亚秒级平均成本达到充分AP和宽带SBL的全部成功水平。候选案例与时间轨迹进一步表明，保留多个可能解并在统一目标下联合比较，可以在有限搜索成本内提高困难构型的恢复能力。波形消元曲率分析则为逐源更新的效率改进提供了依据。后续可将这一框架扩展至未知源数和更一般的传播条件。
